% math4cs-hw5-function-solution.tex

% !TEX program = xelatex
%%%%%%%%%%%%%%%%%%%%
% see http://mirrors.concertpass.com/tex-archive/macros/latex/contrib/tufte-latex/sample-handout.pdf
% for how to use tufte-handout
\documentclass[a4paper, justified]{tufte-handout}

\input{hw-preamble} % feel free to modify this file if you understand LaTeX well
%%%%%%%%%%%%%%%%%%%%
\title{计算机数学 (5-function: 函数)}
\me{魏恒峰}{hfwei@hnu.edu.cn}{}{}
\date{2026年04月09日 \quad 作业 \\ 2026年05月06日 \quad 答案}
%%%%%%%%%%%%%%%%%%%%
\begin{document}
\maketitle
%%%%%%%%%%%%%%%%%%%%
\noplagiarism % PLEASE DON'T DELETE THIS LINE!
%%%%%%%%%%%%%%%%%%%%
% \begin{abstract}
% \end{abstract}
%%%%%%%%%%%%%%%%%%%%
% \beginrequired
%%%%%%%%%%%%%%%

%%%%%%%%%%%%%%%
\begin{problem}[$f$ 与 $f^{-1}$]
  设 $f: A \to B$ 是函数。请证明:
  \begin{enumerate}[(1)]
    \item $f(A_1 \cup A_2) = f(A_1) \cup f(A_2)$
    \item $f^{-1}(B_1 \setminus B_2) = f^{-1}(B_1) \setminus f^{-1}(B_2)$
    \item $B_0 \supseteq f(f^{-1}(B_0))$
  \end{enumerate}
\end{problem}

\begin{proof}
  \begin{enumerate}[(1)]
    \item 对于任意 $b \in B$,
      \begin{align*}
        &b \in f(A_{1} \cup A_{2}) \\
        \iff & \exists a \in A_{1} \cup A_{2}\; f(a) = b \\
        \iff & (\exists a \in A_{1}\; f(a) = b) \lor (\exists a \in A_{2}\; f(a) = b) \\
        \iff & b \in f(A_{1}) \lor b \in f(A_{2}) \\
        \iff & b \in f(A_{1}) \cup f({A_{2}})
      \end{align*}
    \item 对于任意 $a \in A$,
      \setcounter{equation}{0}
      \begin{align*}
        &a \in f^{-1}(B_{1} \setminus B_{2}) \\
        \iff & f(a) \in B_{1} \setminus B_{2} \\
        \iff & f(a) \in B_{1} \land f(a) \notin B_{2} \\
        \iff & a \in f^{-1}(B_{1}) \land \lnot(a \in f^{-1}(B_{2})) \\
        \iff & a \in f^{-1}(B_{1}) \land a \notin f^{-1}(B_{2}) \\
        \iff & a \in f^{-1}(B_{1}) \setminus f^{-1}(B_{2})
      \end{align*}
    \item 对于任意 $b \in B$,
      \setcounter{equation}{0}
      \begin{align*}
        &b \in f(f^{-1}(B_{0})) \\
        \iff & \exists a \in f^{-1}(B_{0})\; b = f(a) \\
        \iff & \exists a \in A\; f(a) \in B_{0} \land b = f(a) \\
        \implies & b \in B_{0}
      \end{align*}
  \end{enumerate}
\end{proof}
%%%%%%%%%%%%%%%

%%%%%%%%%%%%%%%
\newpage
\begin{problem}[$g \circ f$]
  设 $f: A \to B$ 与 $g: B \to C$ 是函数。
  请证明,
  \begin{enumerate}[(1)]
    \item 如果 $f$ 与 $g$ 是满射, 则 $g \circ f$ 是满射。
    \item 如果 $g \circ f$ 是单射, 则 $f$ 是单射。
  \end{enumerate}
\end{problem}

\begin{proof}
  \begin{enumerate}[(1)]
    \item 任取 $c \in C$。\\
      由于 $g$ 是满射,
      \[
        \exists b \in B\; g(b) = c\text{。}
      \]
      任取 $b_{0}$ 满足 $g(b_{0}) = c$。\\
      由于 $f$ 是满射,
      \[
        \exists a \in A\; f(a) = b_{0}\text{。}
      \]
      任取 $a_{0}$ 满足 $f(a_{0}) = b_{0}$。\\
      则有
      \[
        (g \circ f)(a_{0}) = c\text{。}
      \]
    \item 任取 $a_{1}, a_{2} \in A$。\\
      假设 $f(a_{1}) = f(a_{2})$, 则
      \[
        g(f(a_{1})) = g(f(a_{2})),
      \]
      即
      \[
        (g \circ f)(a_{1}) = (g \circ f)(a_{2})\text{。}
      \]
      由于 $g \circ f$ 是单射,
      \[
        a_{1} = a_{2}\text{。}
      \]
  \end{enumerate}
\end{proof}
%%%%%%%%%%%%%%%

%%%%%%%%%%%%%%%
\newpage
\begin{problem}[反函数]
  设 $f: A \to B$ 与 $g: B \to A$ 是函数。
  请证明,
  \[
    (f \circ g = I_B \land g \circ f = I_A) \to g = f^{-1}\text{。}
  \]
\end{problem}

\begin{proof}
  首先, $f \circ g = I_{B}$ 是满射, 因此 $f$ 是满射。\\
  其次, $g \circ f = I_{A}$ 是单射, 因此 $f$ 是单射。\\
  所以, $f$ 是双射。\\
  又由于 $f \circ g = I_{B}$, 因此 $g = f^{-1}$。
  \marginnote{该证明用到了课件中的两个定理, 请自行找出来。}
\end{proof}
%%%%%%%%%%%%%%%

%%%%%%%%%%%%%%%
\newpage
\begin{problem}[函数与等价关系]
  设 $f: X \to Y$ 是满射。
  定义 $X$ 上的二元关系 $R$ 为 $(x, y) \in R$ 当且仅当 $f(x) = f(y)$。
  请证明,
  \begin{enumerate}[(1)]
    \item $R$ 是 $X$ 上的等价关系。
    \item 定义 $h \subseteq (X/R) \times Y$ 为 $h([x]_{R}) = f(x)$。
      请证明, $h$ 是从商集 $X/R$ 到 $Y$ 的函数, 且是满射。
  \end{enumerate}
\end{problem}

\begin{proof}
    \begin{enumerate}[(1)]
    \item 只需证
      \begin{itemize}
        \item $R$ 是自反的。\\
          任取 $x \in X$,
          \[
            f(x) = f(x) \implies (x, x) \in R\text{。}
          \]
        \item $R$ 是对称的。\\
          任取 $x_{1}, x_{2} \in X$,
          \begin{align*}
            & (x_{1}, x_{2}) \in R \\
            \implies & f(x_{1}) = f(x_{2}) \\
            \implies & f(x_{2}) = f(x_{1}) \\
            \implies & (x_{2}, x_{1}) \in R\text{。}
          \end{align*}
        \item $R$ 是传递的。\\
          任取 $x_{1}, x_{2}, x_{3} \in X$,
          \begin{align*}
            & (x_{1}, x_{2}) \in R \land (x_{2}, x_{3}) \in R \\
            \implies & f(x_{1}) = f(x_{2}) \land f(x_{2}) = f(x_{3}) \\
            \implies & f(x_{1}) = f(x_{3}) \\
            \implies & (x_{1}, x_{3}) \in R\text{。}
          \end{align*}
      \end{itemize}
    \item 要证 $h$ 是函数, 只需证
      \begin{itemize}
        \item $\forall S \in X/R\; \exists y \in Y\; h(S) = y$。\\
          设 $X/R = [x]_{R}$。取 $y = f(x)$ 即可。
        \item $\forall S \in X/R\; \exists y_{1}, y_{2} \in Y\;
          ((h(S) = y_{1} \land h(S) = y_{2}) \implies y_{1} = y_{2})$。\\
          根据 $h$ 的定义,
          \[
            \exists x_{1} \in X\; ([x_{1}]_{R} = S \land f(x_{1}) = y_{1}),
          \]
          \[
            \exists x_{2} \in X\; ([x_{2}]_{R} = S \land f(x_{2}) = y_{2})\text{。}
          \]
          因此,
          \begin{align*}
            & [x_{1}]_{R} = [x_{2}]_{R} \\
            \implies & (x_{1}, x_{2}) \in R \\
            \implies & f(x_{1}) = f(x_{2}) \qquad (\because \text{第一题结论}) \\
            \implies & y_{1} = y_{2}
          \end{align*}
      \end{itemize}
      要证 $h$ 是满射, 只需证 $\forall y \in Y\; \exists S \in X/R\; h(S) = y$。\\
      因为 $f$ 是满射, 对于任意 $y \in Y$, 存在 $x_{0}$, 使得 $f(x_{0}) = y$。
      故, 取 $S = [x_{0}]_{R}$ 即可。
  \end{enumerate}
\end{proof}
%%%%%%%%%%%%%%%

%%%%%%%%%%%%%%%%%%%%
% 如果没有反馈，可以把这部分删掉
% \beginfb

% 你可以写 (也可以发邮件)
% \begin{itemize}
%   \item 对课程及教师的建议与意见
%   \item 教材中不理解的内容
%   \item 希望深入了解的内容
%   \item $\cdots$
% \end{itemize}
%%%%%%%%%%%%%%%%%%%%
\end{document}